% Cubed-sphere atlas schematic for App J (Bayesian inference).
% Uses tikz-3dplot for a proper 3-D perspective view.  The sphere
% surface is rendered as a 3-D shaded ball (\shade[ball color]) in
% tdplot_screen_coords with great-circle / latitude wireframe overlaid;
% front-half arcs solid, back-half arcs dashed.  The highlighted face
% is semi-transparent so the projection rays and intersection markers
% behind it remain visible.  Each projection ray carries a marker on
% the cube-face anchor (smaller) AND on the sphere-image (larger).
% Four curved arcs on the sphere bound the sphere-tile (the gnomonic
% image of the highlighted cube-face's edges).
% The sample vector is drawn in two depth-ordered segments: the
% portion inside the cube is rendered BEFORE the transparent face
% (so the face overlays it), and the portion emerging past the face
% is rendered AFTER (so it sits on top).
%
% The figure idiom (cube-tile / sphere-tile naming, gnomonic
% projection ray rendering) follows Chen 2021 (doi:10.1029/2020MS002280).
% The underlying gnomonic projection is Ronchi, Iacono and Paolucci
% (J. Comp. Phys. 124, 93--114, 1996); the cubed-sphere construction
% adapted to coupling-space sampling follows W. E. V. Barker's
% unreleased PSALTer extension shared with this project.
\documentclass[border=3pt,tikz,svgnames]{standalone}
\usepackage{stix}
\usepackage{amsmath}
\usepackage{bm}
\usepackage{tikz}
\usepackage{tikz-3dplot}
\usepackage{tidal-tikz-styles}

\usetikzlibrary{arrows.meta,calc,3d}
\tikzset{>={Stealth[length=1.6mm,width=1.6mm]}}

\colorlet{FaceFill}{PinkFill}
\colorlet{FaceLine}{PinkLine}
\colorlet{TileFill}{BlueFill}
\colorlet{TileLine}{BlueLine}
\colorlet{SphereLine}{GrayLine}
\colorlet{RayLine}{Black!55!GrayLine}
\colorlet{SampleArrow}{PurpleLine}
\colorlet{AxisLine}{GrayLine}

% Camera at inclination 70 deg (from +z) and azimuth 25 deg.  At this
% angle the +x face is the most head-on visible face; the -y face
% (face 2\downarrow) carries the highlight on the back-left side of
% the cube, visible through the unfilled wireframe.
\tdplotsetmaincoords{70}{28}

\begin{document}
\begin{tikzpicture}[
    tdplot_main_coords,
    line cap=round, line join=round,
    font=\footnotesize,
    scale=2.0
  ]

  % Useful sub-tile geometry
  \pgfmathsetmacro{\ttileLo}{-0.3333}
  \pgfmathsetmacro{\ttileHi}{ 0.3333}

  %------------------------------------------------------------------
  % STAGE 0a — inside-sphere axis segments (origin -> unit-sphere
  %                  intersection).  Drawn early so the sphere
  %                  shading overlays them; dashed + lighter to
  %                  visually subordinate the inside portion.
  %------------------------------------------------------------------
  \draw[AxisLine!55, thin, dashed] (0,0,0) -- (1,0,0);
  \draw[AxisLine!55, thin, dashed] (0,0,0) -- (0,1,0);
  \draw[AxisLine!55, thin, dashed] (0,0,0) -- (0,0,1);

  %------------------------------------------------------------------
  % STAGE 2 — far back cube edges (light wireframe of hidden edges)
  %------------------------------------------------------------------
  \draw[FaceLine!40, thin]
    (-1, 1,-1) -- (1, 1,-1) -- (1, 1, 1) -- (-1, 1, 1) -- cycle;
  \draw[FaceLine!40, thin] (-1,-1,-1) -- (-1, 1,-1);
  \draw[FaceLine!40, thin] ( 1,-1,-1) -- ( 1, 1,-1);
  \draw[FaceLine!40, thin] ( 1,-1, 1) -- ( 1, 1, 1);
  \draw[FaceLine!40, thin] (-1,-1, 1) -- (-1, 1, 1);

  %------------------------------------------------------------------
  % STAGE 3 — sample vector segment 1 (origin -> cube-face crossing
  %                  on the -y face).  Drawn BEFORE the transparent
  %                  face so the face overlays it for the "behind"
  %                  depth cue.
  %------------------------------------------------------------------
  \pgfmathsetmacro{\cx}{0.15}
  \pgfmathsetmacro{\cz}{0.10}
  \pgfmathsetmacro{\cn}{sqrt(\cx*\cx + 1 + \cz*\cz)}
  \pgfmathsetmacro{\cxs}{\cx/\cn}
  \pgfmathsetmacro{\cys}{-1/\cn}
  \pgfmathsetmacro{\czs}{\cz/\cn}
  \pgfmathsetmacro{\rmax}{1.55}
  \draw[SampleArrow, very thick]
    (0,0,0) -- (\cx, -1, \cz);
  \fill[SampleArrow] (\cx, -1, \cz) circle (1.2pt);

  %------------------------------------------------------------------
  % STAGE 4 — sphere surface: 3-D shaded ball base (rendered in
  %                  tdplot_screen_coords so the orthographic sphere
  %                  silhouette appears as a screen-space circle with
  %                  proper ball shading) + silhouette outline
  %------------------------------------------------------------------
  \begin{scope}[tdplot_screen_coords]
    \shade[ball color=GrayFill, opacity=0.40] (0,0) circle (1);
    \draw[SphereLine!50, thin] (0,0) circle (1);
  \end{scope}

  %------------------------------------------------------------------
  % STAGE 6 — gnomonic projection rays through 4 cube-face anchors
  %                  on the -y face.  Each ray carries a marker on
  %                  the sphere image AND a smaller marker on the
  %                  cube-face anchor itself.  Drawn BEFORE the
  %                  transparent -y face fill so both markers are
  %                  visible behind it.
  %------------------------------------------------------------------
  % Slight perturbations from the sub-tile centres so the 4 anchors
  % look like a sampled scatter rather than a regular grid.
  \foreach \fx/\fz in {-0.72/-0.58, 0.61/-0.74,
                       -0.55/ 0.68, 0.70/ 0.62}{
    \pgfmathsetmacro{\nrm}{sqrt(\fx*\fx + 1 + \fz*\fz)}
    \pgfmathsetmacro{\sxv}{\fx/\nrm}
    \pgfmathsetmacro{\syv}{-1/\nrm}
    \pgfmathsetmacro{\szv}{\fz/\nrm}
    \draw[RayLine, thin, dotted] (0,0,0) -- (\sxv,\syv,\szv);
    \draw[RayLine, thin] (\sxv,\syv,\szv) -- (\fx, -1, \fz);
    % Sphere-image marker: white-filled ring (suggests "on surface")
    \draw[RayLine, line width=0.6pt, fill=white]
      (\sxv,\syv,\szv) circle (1.5pt);
    \fill[RayLine!85] (\fx, -1, \fz) circle (0.7pt);
  }

  %------------------------------------------------------------------
  % STAGE 7 — sphere-tile boundary curves.  These are the four edges
  %                  of the highlighted (-y) cube-face projected onto
  %                  the sphere by radial normalisation.  Together
  %                  they outline the sphere-tile, the gnomonic image
  %                  of the cube-tile.
  %------------------------------------------------------------------
  % Edge at x = -1, parameter t = z in [-1, 1]
  \draw[FaceLine, thick]
    plot[domain=-1:1, samples=32, smooth]
      ({-1/sqrt(1+1+\x*\x)}, {-1/sqrt(1+1+\x*\x)}, {\x/sqrt(1+1+\x*\x)});
  % Edge at x = +1, parameter t = z
  \draw[FaceLine, thick]
    plot[domain=-1:1, samples=32, smooth]
      ({1/sqrt(1+1+\x*\x)}, {-1/sqrt(1+1+\x*\x)}, {\x/sqrt(1+1+\x*\x)});
  % Edge at z = -1, parameter s = x
  \draw[FaceLine, thick]
    plot[domain=-1:1, samples=32, smooth]
      ({\x/sqrt(\x*\x+1+1)}, {-1/sqrt(\x*\x+1+1)}, {-1/sqrt(\x*\x+1+1)});
  % Edge at z = +1, parameter s = x
  \draw[FaceLine, thick]
    plot[domain=-1:1, samples=32, smooth]
      ({\x/sqrt(\x*\x+1+1)}, {-1/sqrt(\x*\x+1+1)}, {1/sqrt(\x*\x+1+1)});

  %------------------------------------------------------------------
  % STAGE 7b — dashed lines from each cube-tile corner to its
  %                  sphere-tile corner.  These make the gnomonic
  %                  pairing of cube-tile vertices and sphere-tile
  %                  vertices visually explicit.
  %------------------------------------------------------------------
  \foreach \fx/\fz in {-1/-1, 1/-1, 1/1, -1/1}{
    \pgfmathsetmacro{\nrm}{sqrt(\fx*\fx + 1 + \fz*\fz)}
    \pgfmathsetmacro{\sxv}{\fx/\nrm}
    \pgfmathsetmacro{\syv}{-1/\nrm}
    \pgfmathsetmacro{\szv}{\fz/\nrm}
    \draw[FaceLine!70, thin, dashed]
      (\fx, -1, \fz) -- (\sxv, \syv, \szv);
  }

  %------------------------------------------------------------------
  % STAGE 8 — highlighted -y face (transparent fill, sub-tile grid,
  %                  highlighted central sub-tile)
  %------------------------------------------------------------------
  \fill[FaceFill, opacity=0.45]
    (-1,-1,-1) -- (1,-1,-1) -- (1,-1, 1) -- (-1,-1, 1) -- cycle;

  % 3x3 sub-tile grid lines on the -y face
  \foreach \v in {\ttileLo, \ttileHi}{
    \draw[FaceLine!70, thin] (\v,-1,-1) -- (\v,-1, 1);
    \draw[FaceLine!70, thin] (-1,-1,\v) -- ( 1,-1, \v);
  }

  % Highlighted central sub-tile
  \fill[TileFill, opacity=0.40]
    (\ttileLo,-1, \ttileLo) -- (\ttileHi,-1, \ttileLo) --
    (\ttileHi,-1, \ttileHi) -- (\ttileLo,-1, \ttileHi) -- cycle;
  \draw[TileLine, very thick]
    (\ttileLo,-1, \ttileLo) -- (\ttileHi,-1, \ttileLo) --
    (\ttileHi,-1, \ttileHi) -- (\ttileLo,-1, \ttileHi) -- cycle;

  % Face outline
  \draw[FaceLine, very thick]
    (-1,-1,-1) -- (1,-1,-1) -- (1,-1, 1) -- (-1,-1, 1) -- cycle;

  \node[FaceLine, anchor=south west, font=\scriptsize]
    at (1.02,-1.02, -0.95) {face $2\!\downarrow$};

  %------------------------------------------------------------------
  % STAGE 9 — sample vector segment 2 (cube-face -> sphere -> tip)
  %                  drawn AFTER the transparent face so the arrow
  %                  sits on top of the face for the "emerging" cue
  %------------------------------------------------------------------
  \draw[->, SampleArrow, very thick]
    (\cx, -1, \cz) -- ({\rmax*\cxs}, {\rmax*\cys}, {\rmax*\czs});
  % Sphere-image marker for the sample vector: white-filled ring,
  % matching the projection-ray markers for "on the curved surface"
  \draw[SampleArrow, line width=0.6pt, fill=white]
    (\cxs, \cys, \czs) circle (1.3pt);
  \node[SampleArrow, anchor=west, font=\small]
    at ({\rmax*\cxs + 0.04}, {\rmax*\cys + 0.02}, {\rmax*\czs + 0.05})
    {$\theta = r\,\hat r$};

  %------------------------------------------------------------------
  % STAGE 10 — outside-sphere axis segments + arrowheads + labels.
  %                  Drawn last so the bold arrows sit on top of
  %                  everything else, conveying "exits the sphere".
  %------------------------------------------------------------------
  \draw[->, AxisLine, thick] (1,0,0) -- (1.55,0,0)
    node[pos=1.10, AxisLine] {$c_{1}$};
  \draw[->, AxisLine, thick] (0,1,0) -- (0,1.55,0)
    node[pos=1.10, AxisLine] {$c_{2}$};
  \draw[->, AxisLine, thick] (0,0,1) -- (0,0,1.55)
    node[pos=1.10, AxisLine] {$c_{3}$};

  \fill (0,0,0) circle (1pt);
  \node[anchor=north east, font=\scriptsize] at (0,0,-0.02) {$\mathbf 0$};

\end{tikzpicture}
\end{document}
